\scanpage{15}
\deleted{of finite order in the group under consideration \inserted{\deleted{\unsure{stable|stable}} fixed by $\rho$}, one should therefore have}
\[
\deleted{\displaystyle \varepsilon_0\varepsilon_1\varepsilon_\infty=1}
\]
\deleted{i.e. in terms of the $l_i,\sigma_i$}
\[
\deleted{\displaystyle \unsure{l_1|l_1'}\underbrace{\sigma_0l_0}_{\substack{\inserted{\deleted{l_{\unsure{1}}\sigma_\infty}}\\\deleted{l_1^{-1}l_\infty^{-1}\sigma_0}}}l_\infty\underbrace{\sigma_1l_1}_{l_\infty^{-1}l_0^{-1}}l_0\underbrace{\sigma_\infty l_\infty}_{l_0^{-1}l_1^{-1}}=1}
\]
\margin{\deleted{i.e. $l_\infty^{-1}l_0^{-1}l_1^{-1}=1$\quad OK}}
\deleted{We therefore have}
\[
\deleted{\displaystyle \boxed{\varepsilon_0\varepsilon_1\varepsilon_\infty=1}}
\]
\deleted{one can derive prettier expressions for the $\sigma_i$ from this}
\[
\begin{aligned}
\deleted{\displaystyle \sigma_0}&\deleted{\displaystyle =\varepsilon_\infty^{-1}\varepsilon_1}\\
\deleted{\displaystyle \sigma_1}&\deleted{\displaystyle =\varepsilon_0^{-1}\varepsilon_\infty}\\
\deleted{\displaystyle \sigma_\infty}&\deleted{\displaystyle =\varepsilon_1^{-1}\varepsilon_0}
\end{aligned}
\]
\deleted{thus \unsure{also|we have} the unexpected relation}
\[
\deleted{\displaystyle \sigma_\infty\sigma_1\sigma_0=1}
\]
Probably
\[
\varepsilon_0\varepsilon_1\varepsilon_\infty=\sigma_1
\]
i.e.
\[
\begin{aligned}
\varepsilon_0&=\varepsilon_1\varepsilon_\infty\varepsilon_{\unsure{1}}^{-1}\\
\varepsilon_1&=\varepsilon_\infty\varepsilon_0\varepsilon_\infty^{-1}\\
\varepsilon_\infty&=\varepsilon_0\varepsilon_1\varepsilon_0^{-1}
\end{aligned}
\]
\figure{\begin{tikzpicture}\draw (0,0)--(.15,0);\draw (.25,0)--(.8,0);\end{tikzpicture}}

Critique of the notation and of the approach. The elements $l_i,\sigma_i,\varepsilon_i,\rho\ldots$ are in the discrete groups, which would have to be distinguished from the profinite groups --- moreover for the discrete groups one would have to \deleted{\unsure{use}} \unsure{(and|Cat} $\bar m$ for the
\mspage{14}
